\begin{env}[4.2.2]
\label{0.4.2.2-fr}
Sean $\mathcal{M}$, $\mathcal{N}$ dos $\mathcal{A}$-Módulos. Para todo
abierto $U$ de $Y$, se tiene una aplicación canónica
\[
  \Gamma(\psi^{-1}(U),\mathcal{M})
  \times\Gamma(\psi^{-1}(U),\mathcal{N})
  \longrightarrow
  \Gamma(\psi^{-1}(U),\mathcal{M}\otimes_{\mathcal{A}}\mathcal{N})
\]
que es bilineal sobre el anillo
$\Gamma(\psi^{-1}(U),\mathcal{A})=\Gamma(U,\psi_*(\mathcal{A}))$ y,
\emph{a fortiori}, sobre $\Gamma(U,\mathcal{B})$. Define, por tanto,
un homomorfismo
\[
  \Gamma(U,\Psi_*(\mathcal{M}))
  \otimes_{\Gamma(U,\mathcal{B})}
  \Gamma(U,\Psi_*(\mathcal{N}))
  \longrightarrow
  \Gamma(U,\Psi_*(\mathcal{M}\otimes_{\mathcal{A}}\mathcal{N}))
\]
y, como se comprueba de inmediato que estos homomorfismos son compatibles
con las operaciones de restricción, dan un homomorfismo canónico funtorial
de $\mathcal{B}$-Módulos
\[
  \Psi_*(\mathcal{M})\otimes_{\mathcal{B}}\Psi_*(\mathcal{N})
  \longrightarrow
  \Psi_*(\mathcal{M}\otimes_{\mathcal{A}}\mathcal{N})
  \tag{4.2.2.1}\label{0.4.2.2.1-fr}
\]
\oldpage[0\textsubscript{I}]{40}
que, en general, no será ni inyectivo ni sobreyectivo. Si $\mathcal{P}$ es
un tercer $\mathcal{A}$-Módulo, se comprueba de inmediato que el diagrama
\[
\begin{tikzcd}[column sep=large,row sep=large]
\Psi_*(\mathcal{M})\otimes_{\mathcal{B}}\Psi_*(\mathcal{N})
  \otimes_{\mathcal{B}}\Psi_*(\mathcal{P})
  \arrow[r]
  \arrow[d]
&
\Psi_*(\mathcal{M}\otimes_{\mathcal{A}}\mathcal{N})
  \otimes_{\mathcal{B}}\Psi_*(\mathcal{P})
  \arrow[d]
\\
\Psi_*(\mathcal{M})\otimes_{\mathcal{B}}
  \Psi_*(\mathcal{N}\otimes_{\mathcal{A}}\mathcal{P})
  \arrow[r]
&
\Psi_*(\mathcal{M}\otimes_{\mathcal{A}}\mathcal{N}
  \otimes_{\mathcal{A}}\mathcal{P})
\end{tikzcd}
\tag{4.2.2.2}\label{0.4.2.2.2-fr}
\]
es conmutativo.
\end{env}
